%======================================================================
\section{Introduction}
%======================================================================

Graph-based recommenders represent users and items as nodes and logged
interactions as edges, and they score a user--item pair by aggregating signals
along multi-hop walks~\citep{he2020lightgcn,shen2021gfcf}. The logged graph is
not the preference graph. An edge exists only if the platform \emph{exposed}
the item and the user then responded, so the graph records the exposure policy
as much as the user's taste. Causal learning offers two families of
corrections. \emph{Loss-level} debiasing re-weights the training objective by
inverse propensities or doubly robust (DR)
terms~\citep{schnabel2016,wang2019drjl,saito2020unbiased}, but it leaves the
propagation operator of a graph model untouched. \emph{Propagation-level}
debiasing changes the propagated graph, for example by weighting each edge by
its inverse propensity~\citep{kim2022navip} or by learning popularity-aware
aggregation weights~\citep{zhou2023apda,dpaa2026}.

This paper starts from a simple observation about the second family. An
inverse-propensity edge estimate $W_e=O_eY_e/p_e$ is unbiased for the potential
outcome $Y_e$, where $O_e$ indicates that edge $e$ was exposed, $Y_e$ is the
outcome and $p_e$ the exposure propensity. A walk that traverses $e$ twice
contributes $W_e^2$, and $\E[W_e^2]=Y_e/p_e\neq Y_e$. The same holds for DR edge
estimates. Walks that repeat an edge are part of every multi-hop operator: the
back-and-forth step $u\to j\to u$ of a three-hop walk is one. Their excess is
largest where the propensity is smallest, that is, on the rarely exposed edges
that debiasing is meant to recover.

%======================================================================
\section{Research objectives and contributions}\label{sec:objectives}
%======================================================================

The aim of this study is to determine whether and when removing the
repeated-walk bias helps graph recommendation from exposure-biased logs. We
pursue four objectives, each tied to a research question (RQ).

\begin{description}
\item[O1 (estimator).] Derive the exact bias of inverse-propensity and DR graph
propagation and an estimator that removes it. \emph{RQ1: how large is the bias,
and can it be removed exactly at a cost comparable to the uncorrected operator?}
\item[O2 (conditions).] State the conditions under which the correction is
unbiased and test them in the way practitioners fit nuisances. \emph{RQ2: do the
guarantees survive cross-fitting and data-dependent degree weights?}
\item[O3 (value).] Measure what the correction changes. \emph{RQ3: does it
improve the accuracy of recommendations against tuned baselines on four
benchmarks, and does it improve the values of the scores on a log with known
exposure?}
\item[O4 (invariance).] Test on real logs the prediction that expected scores do
not depend on an item's exposure. \emph{RQ4: how informative is a thinning audit
of this property?}
\end{description}

\paragraph{Answers in brief.} The bias is exact and removable at three hops for
the cost of one item Gram matrix (RQ1). The guarantees need nuisances that do
not depend on the log: cross-fitting suffices if the degree weights are fixed a
priori, but not if they are computed from the log, which our experiments do
(RQ2). The correction leaves accuracy unchanged on four datasets, and it
improves score values only under oracle conditions: true propensities and
weights fixed in advance (RQ3). The thinning audit has little power and the
selected configurations are not invariant on two of four datasets (RQ4).

\paragraph{Contributions.}
\begin{enumerate}
\item \emph{Diagnosis.} The exact conditional bias of inverse-propensity and DR
propagation on unexposed candidates; under a propensity clip $\tau$ its worst
case is $\Theta(1/\tau)$ (Theorem~\ref{thm:lower}).
\item \emph{Estimator.} DRUP (Doubly Robust, walk-Unbiased Propagation) uses
$Y_e^2=Y_e$ for binary outcomes and replaces each repeated edge by its
idempotent unbiased counterpart. The three-hop operator has the asymptotic cost
of the uncorrected one, and Möbius inversion over the coincidence patterns of
walk indices makes the correction exact for any odd number of hops
(Theorem~\ref{thm:khop}). A sample split supplies nuisances independent of the
edge estimates (DRUP-split).
\item \emph{Guarantees and their scope.} With nuisances independent of the log,
DRUP is edge-wise doubly robust for propagation over the full-exposure graph
(Theorem~\ref{thm:unbiased}), unbiased on unexposed candidates up to the
imputation error at the candidate (Corollary~\ref{cor:cand}), and its expected
scores do not respond to interventions on an item's exposure
(Theorem~\ref{thm:fair}). We show by simulation where these guarantees fail.
\item \emph{Evidence.} A semi-synthetic test on KuaiRec with known exposure;
a comparison with up to 26 tuned baselines on Coat, Yahoo!\,R3, KuaiRand and
KuaiRec with user-level tests; and real-data thinning audits with a negative
control. The properties that follow from linearity alone (attributions,
certificates, private release, capped allocation) are in the Supplementary
Material.
\end{enumerate}

\paragraph{Outline.} Section~\ref{sec:related} reviews related work,
Section~\ref{sec:drup} presents the method, Section~\ref{sec:setup} the datasets,
baselines and protocol, Section~\ref{sec:results} the results, and
Section~\ref{sec:disc} the theoretical and practical implications and limits.

%======================================================================
\section{Related work}\label{sec:related}
%======================================================================

\paragraph{Graph-based recommendation.}
Graph neural recommenders propagate embeddings over the user--item graph.
LightGCN~\citep{he2020lightgcn} removes feature transformations and
non-linearities, which reduces propagation to a weighted sum of powers of a
normalised adjacency; UltraGCN~\citep{mao2021ultragcn} approximates
infinite-layer propagation, and SGL~\citep{wu2021sgl} and
SimGCL~\citep{yu2022simgcl} add contrastive views. Closed-form models are strong
competitors: EASE~\citep{steck2019ease} and EDLAE~\citep{steck2020edlae} learn an
item--item matrix with a constrained diagonal, which, like our diagonal
correction, removes an item's own contribution to its score;
GF-CF~\citep{shen2021gfcf} and BSPM~\citep{choi2023bspm} cast collaborative
filtering as graph signal filtering. Carefully tuned simple baselines often
match graph models~\citep{anelli2023myth}. All of these propagate over the
\emph{logged} graph. DRUP keeps the linear propagation form but propagates over
an unbiased estimate of the full-exposure graph.

\paragraph{Debiasing in recommendation.}
Logs are missing not at random~\citep{marlin2009collaborative}; surveys organise
the resulting biases and cast debiasing as causal
inference~\citep{chen2023bias,gao2024causalsurvey}. Propensity-based methods
re-weight the loss~\citep{schnabel2016,swaminathan2015snips,saito2020unbiased}.
Doubly robust learning adds an imputation model~\citep{wang2019drjl}, and later
variants reduce its bias and variance~\citep{guo2021mrdr,dai2022gdr,li2023stabledr}.
A small uniform sample can supervise the debiasing~\citep{chen2021autodebias}.
Causal approaches to popularity bias use counterfactual
inference~\citep{wei2021macr}, confounder adjustment~\citep{zhang2021pda} or
disentangled representations~\citep{zheng2021dice}; recent work in this journal
studies popularity debiasing from both sides~\citep{boratto2021connecting}, with a
controllable trade-off~\citep{okamura2024flexibly}, through invariant
learning~\citep{bai2025invariant} and through causal deconfounding of
multi-criteria ratings~\citep{guo2026confounding}. These methods correct the loss
or the output of a model; DRUP applies inverse-propensity and DR estimation to the
\emph{edges} of the propagated graph.

\paragraph{Debiasing inside graph propagation.}
NAVIP~\citep{kim2022navip} weights each edge by its inverse propensity inside
LightGCN. r-AdjNorm~\citep{zhao2022radjnorm} tunes the adjacency normalisation
that controls the amplification of popularity, and graph convolutions provably
amplify popularity bias~\citep{chen2024gcnpop}. APDA~\citep{zhou2023apda},
CAGED~\citep{que2025caged} and DPAA~\citep{dpaa2026} learn aggregation weights.
These methods change \emph{which} adjacency is propagated. None analyses what
happens to an unbiased edge estimate when the re-weighted adjacency is multiplied
with itself over several hops. Theorem~\ref{thm:lower} shows that the product is
then biased by a term of order $1/p$ that concentrates on rarely exposed items.
Our linear baseline over the logged graph searches the normalisation exponent and
is therefore a training-free r-AdjNorm; we also train r-AdjNorm and NAVIP.

\paragraph{Unbiased estimation of polynomials under missingness.}
An unbiased estimator of a polynomial in unknown quantities must avoid products of
the same noisy measurement, as in U-statistics~\citep{hoeffding1948}. For
Bernoulli-masked data the masked Gram matrix needs a $1/p^2$ correction off the
diagonal and $1/p$ on it~\citep{lounici2014}; our item operator is the analogue
with heterogeneous propensities. Möbius inversion on the partition
lattice~\citep{rota1964} turns homomorphism counts into injective counts and is
the combinatorial counterpart of our $K$-hop construction. Non-backtracking
operators~\citep{krzakala2013} and self-avoiding walks~\citep{hopkins2017bayesian}
discard repeated walks; our correction keeps them and replaces each repeated
factor by its unbiased counterpart. In recommendation,
\citet{schnabel2020item} debias item co-occurrence with propensities. The
reduction itself is therefore not new; what is new is its combination with
inverse-propensity and DR edge estimates and degree weights, the exact bias of the
uncorrected operators, and the consequences for exposure invariance. Our
treatment of nuisances follows double machine learning~\citep{chernozhukov2018dml}.

\paragraph{Fairness, explanation, robustness and privacy.}
Fairness notions cover users and items, groups and individuals, and outcomes and
exposure~\citep{wang2023fairness,amigo2023fairness}; item-side exposure can be
allocated in proportion to merit~\citep{singh2018exposure,diaz2020expected}, and
multi-stakeholder views weigh users, providers and the
platform~\citep{abdollahpouri2020multistakeholder,ekstrand2022fairness}. DRUP
contributes a causal property of the merit \emph{estimate}, exposure invariance,
which is not a fairness criterion. Graph explainers search for important or
minimal edge sets~\citep{ying2019gnnexplainer,lucic2022cfgnn,ghazimatin2020prince},
recommenders can be attacked by injected
profiles~\citep{lam2004shilling,li2016poisoning,fang2018poisoning}, and private
recommenders perturb statistics or matrix
completion~\citep{mcsherry2009netflix,jain2018dpmc}. Any linear operator with fixed
nuisances admits exact removal attributions, closed-form certificates and a jointly
private release of its item operator~\citep{balle2018analytic,hsu2014matchings};
these are in the Supplementary Material.

\begin{table}[t]\centering\footnotesize
\caption{DRUP and representative related methods. ``loss'': the correction
acts on the training objective; ``output'': it acts on the predictions;
``learned'': learned aggregation weights; ``--'': not applicable. Properties
that follow from linearity (exact attributions, certificates, joint DP) hold
for every linear operator, including linear LightGCN, GF-CF and DR adjacency,
and are not listed.}\label{tab:related}
\setlength{\tabcolsep}{4pt}
\begin{adjustbox}{max width=\textwidth}
\begin{tabular}{lcccc}
\toprule
Method & \shortstack{debiases\\propagation} & \shortstack{unbiased\\beyond 1 hop} & \shortstack{doubly\\robust} &
\shortstack{exposure\\invariance}\\
\midrule
\shortstack[l]{LightGCN, NGCF,\\UltraGCN, SGL, SimGCL} & no & no & no & no\\[2pt]
\shortstack[l]{linear LightGCN,\\EASE, GF-CF, BSPM} & no & no & no & no\\[2pt]
\shortstack[l]{IPS, SNIPS, DR,\\MRDR, StableDR} & loss & -- & DR variants & no\\[2pt]
MACR, PDA, DICE & output & -- & no & no\\
NAVIP & edge weights & no & no & no\\
\shortstack[l]{r-AdjNorm, APDA,\\CAGED, DPAA} & learned & no & no & no\\[2pt]
DR adjacency (uncorrected) & DR edges & no & edge-wise & no$^\ast$\\
\midrule
DRUP-split (ours; assumptions hold) & DR edges & any odd $K$ & yes$^\dagger$ & in expectation$^\dagger$\\
DRUP, cross-fitted (ours; as run) & DR edges & no$^\ddagger$ & no$^\ddagger$ & no$^\ddagger$\\
\bottomrule
\end{tabular}
\end{adjustbox}
\par\smallskip
\raggedright\scriptsize $^\ast$Over-corrects (Theorem~\ref{thm:fair}) with nuisances independent of the log; as invariant as DRUP in our real-data audits, which cannot tell them apart.
$^\dagger$Guaranteed under Assumptions~\ref{as:unconf}--\ref{as:fixed} (the sample split) and unclipped propensities.
$^\ddagger$No guarantee: nuisances are cross-fitted on the same log and degree weights and score constants are computed from it
(Remark~\ref{rem:xfit}); in simulation the correction then no longer reduces bias, and the audited configurations are not invariant on Yahoo!\,R3 and KuaiRand.
The cross-fitted operator is what the accuracy tables report; its accuracy equals that of uncorrected DR adjacency.
\end{table}
